\documentclass[11pt]{article}
\usepackage{amsmath,amssymb,amsthm,mathtools}
\usepackage[margin=1in]{geometry}
\usepackage{booktabs,array,longtable}
\usepackage{hyperref}
\setlength{\emergencystretch}{3em}

\newtheorem{principle}{Principle}
\newtheorem{conjecture}{Conjecture}
\newtheorem{theorem}{Theorem}
\newtheorem{proposition}{Proposition}
\newtheorem{lemma}{Lemma}
\newtheorem{remark}{Remark}

\title{Amplification Must Be Paid:\\
A Conditional Source-Ledger Framework for Navier--Stokes}
\author{Joseph Roy}
\date{June 2026}

\begin{document}
\maketitle

\begin{abstract}
I propose a scale-critical accounting principle for three-dimensional
Navier--Stokes dynamics: persistent amplification must be paid by an
identifiable channel.  The guiding intuition is that constraints are not
created or destroyed by a blowup scenario; they are redistributed among
source, flux, strain, renewal, quotient, separation, and residual channels.
The proposed ledger tries to make that redistribution visible.

The framework is built around the positive vortex-stretching density
\(\sigma_+=(\omega\cdot S\omega)_+\) and a scale-critical denominator ledger
\(D(r)=\mathbf 1^T X(r)\).  The source ledger is formulated under a local
source-finiteness or truncation hypothesis,
\[
P(r_0)=r_0\iint_{Q_{r_0}(z_0)}\sigma_+\,dx\,dt<\infty,
\]
which is automatic in smooth pre-singular regions but is not automatic in the
bare Leray--Hopf energy class.

This manuscript does not claim a complete proof of Navier--Stokes regularity.
It gives a stand-alone statement of the amplification-payment program, defines
the live ledger rows, records partial closure mechanisms, states the conditional
regularity theorem with its constant gap, reports adversarial numerical tests
only as mechanism tests, and lists the remaining theorem gates.
\end{abstract}

\clearpage
\section{Introduction}
\label{sec:introduction}

The Navier--Stokes equations have conservation and dissipation structures, but
possible singularity formation is a problem of amplification.  The guiding idea
of this manuscript is that amplification should not be free.  If vorticity,
strain, boundary flux, annular leakage, or coherent residuals persist at a
scale-critical level, then the dynamics should generate a visible payment in a
finite ledger.

The proposed law is not a new PDE.  It is a derived accounting principle:
\[
\boxed{\text{amplification must be paid}.}
\]
This paper is designed to stand alone as the mechanism paper: it defines the
ledger, states what is proved conditionally, explains what has been tested, and
makes the remaining proof gates explicit.  The downstream technical gates K3
through K7 are isolated in the companion note
\texttt{technical\_gates\_k3\_k7.pdf}.

\clearpage
\section{The Proposed Law}
\label{sec:law}

The intuition behind the proposed law is simple: constraints cannot be created
or destroyed for free; they can only be redistributed.  If this intuition is
right, then a Navier--Stokes growth scenario cannot create persistent
scale-critical amplification without sending the cost somewhere visible.  The
ledger below is one way to make that intuition falsifiable: amplification must
route to source, boundary, tail, renewal, quotient redundancy, separation,
coherent residual, or an explicitly unresolved row.

\begin{principle}[Amplification payment principle]
\label{prin:payment}
Persistent scale-critical amplification in three-dimensional incompressible
Navier--Stokes should generate an identifiable payment channel.  In the
source-ledger formulation, this takes the schematic form
\begin{equation}
\label{eq:payment-schematic}
\mathrm{Amplification}(r)\le \mathrm{Payment}(r)+o_r(1),
\end{equation}
where \(\mathrm{Payment}(r)\) is a finite ledger of source, boundary, tail,
strain, renewal, quotient, separation, and residual charges.  In the
source-ledger formulation below, the amplification and payment terms are
represented through \(P(r)\), \(D(r)\), and the final physical ledger.
\end{principle}

\begin{remark}[Conjectural status]
\label{rem:conjectural-status}
Principle~\ref{prin:payment} is a conjectural organizing law.  The paper does
not assert \eqref{eq:payment-schematic} as an unconditional theorem.  The
conditional theorem in Section~\ref{sec:conditional-theorem} states exactly
which remaining gates would turn the law into a regularity contradiction.
\end{remark}

Throughout, \(o_r(1)\) denotes a quantity tending to zero as \(r\downarrow0\)
along the stated sequence of scales, with all fixed ledger parameters held
constant.  When the proof uses selected scales \(r_n\downarrow0\), I also
write \(o_n(1)\).

\subsection*{Standalone versus technical-ledger dependencies}

This manuscript is designed to be readable without the full technical ledger,
but it does not make every ledger estimate self-contained.  It proves or
records the following ingredients:
\begin{enumerate}
\item the conditional amplification-payment closure theorem;
\item the source-doubling scale-selection lemma;
\item fixed-dilation and bounded-macroscopic-parent tail payment;
\item the empirical attack summaries as mechanism tests.
\end{enumerate}
It imports the following technical-ledger inputs:
\begin{enumerate}
\item denominator-to-CKN compatibility;
\item K3 quotient-or-charge for the final physical ledger;
\item K4 forward absorption of the final physical ledger;
\item K5 reverse matrix constants, including packed boundary-row constants;
\item K5B high-strain transport rigidity;
\item K6 spectral-radius and loop-gap certification.
\end{enumerate}
Thus the paper should be read as a conditional program note: it isolates the
finite list of theorem gates rather than claiming that all gates are already
proved here.

\clearpage
\section{Source Ledger and Denominator}
\label{sec:denominator}

Fix a candidate point \(z_0=(x_0,T)\) and set
\[
Q_r(z_0):=B_r(x_0)\times(T-r^2,T).
\]
Let \(S=(\nabla u+\nabla u^T)/2\), \(\omega=\nabla\times u\), and
\[
\sigma_+:=(\omega\cdot S\omega)_+.
\]
The source ledger is
\begin{equation}
\label{eq:P-def}
P(r):=r\iint_{Q_r(z_0)}\sigma_+\,dx\,dt.
\end{equation}

\begin{remark}[Source-finiteness input]
\label{rem:source-finiteness}
For fixed parent cylinders used in the ledger, I assume or import
\[
P(r_0)=r_0\iint_{Q_{r_0}(z_0)}\sigma_+\,dx\,dt<\infty.
\]
In smooth pre-singular regions this follows from local smoothness.  For
suitable weak solutions it is a serious additional source-ledger hypothesis,
or else it must be replaced by a truncation/renormalization version of the
framework.
\end{remark}

Let \(\phi_r\) be a standard nonnegative cutoff adapted to \(Q_{2r}(z_0)\),
equal to one on \(Q_r(z_0)\), with
\[
|\nabla\phi_r|\lesssim r^{-1},\qquad
|\partial_t\phi_r|+\nu|\Delta\phi_r|\lesssim r^{-2}.
\]
The live denominator is
\begin{equation}
\label{eq:D-def}
D(r):=\mathbf 1^T X(r),
\qquad
X(r):=
\begin{pmatrix}
\mathcal E_\omega(r)\\
\mathcal B_\omega(r)\\
\mathcal D_{\mathrm{ARR}}(r)\\
\mathfrak B_{\mathrm{strain}}(r)\\
M_{\mathrm{deact}}(r)\\
\mathcal E_\nu^{\mathrm{coh},+}(r)
\end{pmatrix}.
\end{equation}
All rows of \(D(r)\) are required to be nonnegative, scale-critical, and stable
under the scale subsequences used in Theorem~\ref{thm:conditional-closure}.  The
normalizations used here are the following.

\begin{enumerate}
\item \textbf{Vorticity energy/dissipation.}
\[
\begin{aligned}
\mathcal E_\omega(r)
&:= r\operatorname*{ess\,sup}_{t\in(T-4r^2,T)}
      \int_{B_{2r}}|\omega|^2\phi_r^2\,dx \\
&\quad+\nu r\iint_{Q_{2r}}|\nabla\omega|^2\phi_r^2\,dx\,dt .
\end{aligned}
\]

\item \textbf{Boundary/collar flux.}
\[
\mathcal B_\omega(r)
:=\mathcal B_\omega^{\mathrm{sh}}(r)
  +\mathcal B_{\omega,\lambda_0}^{\mathrm{pack}}(2r),
\]
where
\[
\mathcal B_\omega^{\mathrm{sh}}(r)
:=r\iint_{Q_{2r}\setminus Q_r}
|\omega|^2\Theta_{\phi_r}\,dx\,dt
\]
and
\[
\Theta_{\phi_r}
:=|\partial_t\phi_r|+\nu|\Delta\phi_r|
  +|u-\bar u_{B_{2r}}(t)|\,|\nabla\phi_r|.
\]

\item \textbf{Annular renewal/tail reservoir.}
\[
\mathcal D_{\mathrm{ARR}}(r)
:=S_{\mathrm{ARR}}^{\mathrm{tail}}(r)+o_r(1),
\]
with
\[
S_{\mathrm{ARR}}^{\mathrm{tail}}(r)
:=\sum_{k\ge0}2^{-k}\,2^{k+1}r
\iint_{\mathcal A_k(r)}
\sigma_+\chi_k^{\mathrm{tail}}\,dx\,dt .
\]
Here
\[
\mathcal A_k(r)
:=Q_{2^{k+1}r}(z_0)\setminus Q_{2^kr}(z_0)
\]
inside the chosen parent cylinder.

\item \textbf{Source-weighted strain demand.}
\[
\mathfrak B_{\mathrm{strain}}(r)
:=r\iint_{\mathcal U_r}\sigma_+S_r\phi_r^2\,dx\,dt,
\qquad
S_r:=r^2|S_{\mathrm{loc}}|.
\]
The raw term \(r\iint\sigma_+|S_{\mathrm{loc}}|\phi_r^2\) is excluded as
supercritical.

\item \textbf{Source deactivation.}
\[
M_{\mathrm{deact}}(r)
:=r\iint_{\cup_j\mathcal T_{r,j}\cap Q_r}
\sigma_+\mathbf 1_{\mathrm{deact}}\phi_r^2\,dx\,dt+o_r(1),
\]
where \(\mathbf 1_{\mathrm{deact}}\) marks promoted source that fails the
persistence/deactivation test.

\item \textbf{Coherent viscous residual.}
\[
\mathcal E_\nu^{\mathrm{coh},+}(r)
:=r\iint_{Q_{2r}}\mathcal R_{\nu,\mathrm{coh}}^+\,dx\,dt.
\]
This is the positive coherent part of the viscous/projection residual left by
the localized vorticity ledger after source, boundary, and tail rows are
separated.
\end{enumerate}

\subsection*{Technical-ledger symbols}

The following symbols are not new proof inputs in this paper; they are the
named placeholders used by the technical ledger.
\begin{itemize}
\item \(S_{\mathrm{loc}}\) is the locally normalized strain tensor used in the
source-weighted strain row, after subtracting the harmless large-scale or
harmonic component fixed by the cutoff convention.
\item \(\mathcal T_{r,j}\) are promoted finite-thickness source tubes generated
by the primary-source promotion theorem.
\item \(\chi_k^{\mathrm{tail}}\) is the tail-membership cutoff assigning
promoted source mass in the \(k\)-th annular parent shell to the ARR row.
\item \(\mathbf 1_{\mathrm{deact}}\) marks promoted source packets that fail the
persistence/renewal test and are therefore routed to \(M_{\mathrm{deact}}\).
\item \(\mathcal R_{\nu,\mathrm{coh}}^+\) is the positive coherent part of the
viscous/projection residual after source, boundary, tail, and strain terms have
been separated in the localized vorticity ledger.
\item Quotient redundancy means that overlapping promoted packets are
identified by chains of significant-overlap pairs whose final physical distance
is at most the quotient threshold \(\delta\).
\end{itemize}
The companion technical ledger must provide exact constructions and constants
for these objects.  This paper records where each object enters the conditional
closure.

The packed boundary contribution is
\begin{equation}
\label{eq:packed-boundary}
\mathcal B_{\omega,\lambda_0}^{\mathrm{pack}}(R)
:=
\sup_{\substack{0<a<b\\ \lambda_0 b\le R}}
\int_a^b
\mathcal B_{\omega}^{\mathrm{adv}}(s,\lambda_0)
\frac{ds}{s},
\end{equation}
where
\[
\mathcal B_{\omega}^{\mathrm{adv}}(s,\lambda_0)
:=
s\iint_{Q_{\lambda_0s}\setminus Q_s}
|u-\bar u_{B_{\lambda_0s}}(t)|\,|\omega|^2\,
|\nabla\phi_{s,\lambda_0}|\phi_{s,\lambda_0}\,dx\,dt.
\]
This packed row is essential for the good-scale null-source result; replacing
it by a single shell reopens the boundary contraction problem.

\clearpage
\section{Physical Charges and Quotient Ledger}
\label{sec:physical-ledger}

The final physical ledger is schematically
\begin{equation}
\label{eq:Fphys-final}
\mathcal F_{\mathrm{phys}}^{\mathrm{final}}
:=
\mathcal F_{\mathrm{phys}}
+R_{\mathrm{osc}}^{\mathrm{fam}},
\end{equation}
where \(\mathcal F_{\mathrm{phys}}\) includes source-weighted shape,
renewal, spread/separation, activity, boundary, tail, and coherent-residual
coordinates.  The oscillation/twist coordinate is explicit because large
accumulated strain need not imply large net shape growth.

For promoted branches \(i,j\), the final physical distance has the form
\begin{equation}
\label{eq:dphys-final}
d_{\mathrm{phys}}^{\mathrm{final}}(i,j)
:=
\max\{d_{\mathrm{phys}}(i,j),w_{\mathrm{osc}}\Phi_{\mathrm{osc}}(i,j)\}.
\end{equation}
The quotient relation \(i\sim_\delta j\) identifies branches connected by
significant-overlap chains with
\[
d_{\mathrm{phys}}^{\mathrm{final}}\le\delta.
\]
The nonredundant edge set is
\[
E_{\mathrm{nonred}}^{\mathrm{final}}
:=\{(i,j):i\not\sim_\delta j,\ m_{ij}>0\},
\]
where \(m_{ij}\) is the source-overlap edge mass.  The K3 target is
\begin{align}
\label{eq:K3-target-standalone}
\sum_{(i,j)\in E_{\mathrm{nonred}}^{\mathrm{final}}}m_{ij}
&\ge c_m^{\mathrm{final}}\eta P(r),\\
\sum_{(i,j)\in E_{\mathrm{nonred}}^{\mathrm{final}}}
m_{ij}d_{\mathrm{phys}}^{\mathrm{final}}(i,j)
&\le C_*^{\mathrm{final}}P(r)
\mathcal F_{\mathrm{phys}}^{\mathrm{final}}(r).
\end{align}
Together these imply
\begin{equation}
\label{eq:cfinal-def}
\mathcal F_{\mathrm{phys}}^{\mathrm{final}}(r)
\ge
c_{\mathrm{final}}(\eta),
\qquad
c_{\mathrm{final}}(\eta):=
\frac{\delta c_m^{\mathrm{final}}}{C_*^{\mathrm{final}}}\eta.
\end{equation}

\clearpage
\section{Regularity Bridge}
\label{sec:ckn-bridge}

The final regularity step is imported from standard local regularity theory,
but the normalization must be explicit.  Define the mean-subtracted CKN
quantity
\begin{equation}
\label{eq:KCKNms}
\mathcal K_{\mathrm{CKN}}^{\mathrm{ms}}(r;z_0)
:=
r^{-2}\iint_{Q_r(z_0)}
\Big(|u-\bar u_{B_r}(t)|^3
+|p-\bar p_{B_r}(t)|^{3/2}\Big)\,dx\,dt.
\end{equation}

\begin{proposition}[Imported denominator-to-CKN bridge]
\label{prop:CKN-bridge}
Assume the denominator compatibility used in the technical ledger:
\[
D(r_n)\to0
\quad\Longrightarrow\quad
r_n\iint_{Q_{2r_n}}|\omega|^3\,dx\,dt\to0
\]
and the harmless cutoff, harmonic, and pressure-projection remainders vanish.
Then
\[
D(r_n)\to0
\quad\Longrightarrow\quad
\mathcal K_{\mathrm{CKN}}^{\mathrm{ms}}(r_n;z_0)\to0.
\]
Consequently, for sufficiently large \(n\), the CKN epsilon-regularity
criterion applies at scale \(r_n\), so \(z_0\) is regular.
\end{proposition}

\begin{proof}[Imported mechanism]
The bridge uses local Gagliardo--Nirenberg interpolation for vorticity, local
div--curl/Biot--Savart control of the mean-subtracted velocity, local pressure
decomposition with Calderon--Zygmund estimates, and the mean-subtracted
Caffarelli--Kohn--Nirenberg epsilon criterion.  No source-packing argument is
used in this final imported step.
\end{proof}

\clearpage
\section{Conditional Closure Algebra}
\label{sec:conditional-theorem}

The theorem below is an algebraic closure theorem.  Its role is to show that
the listed analytic inputs are sufficient for regularity.  It does not prove
the imported K3--K7 inputs.

\begin{theorem}[Algebraic conditional amplification-payment closure]
\label{thm:conditional-closure}
Let \(z_0\) be a candidate singular point.  Suppose the following inputs hold
along a scale sequence \(r_n\downarrow0\):
\begin{enumerate}
\item \textbf{Persistent source and promotion:}
\[
P(r_n)\ge p_*>0
\]
and a finite-thickness promoted source family captures a fixed fraction
\(\eta>0\).
\item \textbf{Quotient-or-charge:} either the promoted family is in the
nonredundant charge branch and
\[
\mathcal F_{\mathrm{phys}}^{\mathrm{final}}(r_n)
\ge c_{\mathrm{final}}(\eta)
\]
holds, or the quotient-redundant branch has already been routed to a
null-source/good-scale alternative that gives regularity.
\item \textbf{Forward absorption:}
\[
P(r_n)\mathcal F_{\mathrm{phys}}^{\mathrm{final}}(r_n)
\le
C_{\mathrm{forw}}^{\mathrm{final}}D(r_n)+o_n(1).
\]
\item \textbf{Reverse matrix absorption:}
\[
X(r_n)\le AX(r_n)+bP(r_n)+o_n(1),
\]
where \(A\ge0\), \(b\ge0\), and \(\rho(A)<1\).  Here \(o_n(1)\) in the vector
inequality denotes a vector whose components tend to zero; inequalities are
interpreted componentwise.  Set
\[
C_0:=\mathbf 1^T(I-A)^{-1}b.
\]
\item \textbf{Strict loop gap:}
\[
C_{\mathrm{forw}}^{\mathrm{final}}C_0
<
c_{\mathrm{final}}(\eta).
\]
\item \textbf{CKN bridge:} Proposition~\ref{prop:CKN-bridge} applies to the
chosen denominator \(D\).
\end{enumerate}
Then \(z_0\) is regular.
\end{theorem}

\begin{proof}
If the quotient-redundant alternative has already been routed to a
null-source/good-scale regularity conclusion, there is nothing further to
prove.  Otherwise the contradiction sequence is in the nonredundant charge
branch.  By
quotient-or-charge and forward absorption,
\[
c_{\mathrm{final}}(\eta)P(r_n)
\le
P(r_n)\mathcal F_{\mathrm{phys}}^{\mathrm{final}}(r_n)
\le
C_{\mathrm{forw}}^{\mathrm{final}}D(r_n)+o_n(1).
\]
Since \(\rho(A)<1\), the reverse matrix inequality gives
\[
X(r_n)\le (I-A)^{-1}bP(r_n)+o_n(1),
\]
and therefore
\[
D(r_n)=\mathbf 1^TX(r_n)
\le
C_0P(r_n)+o_n(1).
\]
Substitution yields
\[
c_{\mathrm{final}}(\eta)P(r_n)
\le
C_{\mathrm{forw}}^{\mathrm{final}}C_0P(r_n)+o_n(1).
\]
Because \(C_{\mathrm{forw}}^{\mathrm{final}}C_0<c_{\mathrm{final}}(\eta)\)
and \(P(r_n)\ge p_*>0\), this is impossible for large \(n\).  Thus the
persistent-source nonredundant charge branch is ruled out, while the quotient
branch is handled by the routed quotient/null-source alternative.  Equivalently,
the loop forces \(P(r_n)\to0\) and then \(D(r_n)\to0\);
Proposition~\ref{prop:CKN-bridge} gives regularity.
\end{proof}

\begin{remark}[Weighted spectral certificate]
The condition \(\rho(A)<1\) may be certified without eigenvalue calculation.
If \(w_i>0\) and
\[
q_A(w):=\max_i\frac{1}{w_i}\sum_j A_{ij}w_j<1,
\]
then \(\rho(A)\le q_A(w)<1\).  A usable reverse constant is
\[
C_0^{\mathrm{cert}}(w)
:=
\frac{(\sum_iw_i)\max_i(b_i/w_i)}{1-q_A(w)}.
\]
The certified loop gap is
\[
C_{\mathrm{forw}}^{\mathrm{final}}C_0^{\mathrm{cert}}(w)
<c_{\mathrm{final}}(\eta).
\]
\end{remark}

\clearpage
\section{Partial Closure Results}
\label{sec:partial-results}

The following mechanisms are currently the most stable pieces of the program.
They are stated here in simplified form; their detailed proofs and proof
obligations are reflected in the technical ledger.

\begin{proposition}[Good-scale null-source regularity]
\label{prop:null-source}
If the packed boundary row \eqref{eq:packed-boundary} is part of \(D(r)\) and
\[
\limsup_{r\downarrow0}P(r)=0,
\]
then there exists \(r_j\downarrow0\) such that \(D(r_j)\to0\), hence
\(z_0\) is regular by Proposition~\ref{prop:CKN-bridge}.
\end{proposition}

\begin{proof}[Mechanism]
The localized vorticity identity separates the collar flux into a
source-active part and a no-source advective part.  The source-active part is
paid by \(P\).  The no-source part obeys a logarithmic packing estimate:
\[
\int_r^{\Lambda r}
\mathcal B_\omega^{\mathrm{adv,nosrc}}(s,\lambda_0)
\frac{ds}{s}
\le
C_{\mathrm{avg}}D(\Lambda\lambda_0 r)+o_r(1).
\]
Pigeonholing selects \(s\in[r,\Lambda r]\) with contraction coefficient
\[
\theta_*=
\theta_{\mathrm{cut}}(\Lambda,\lambda_0)
+\frac{C_{\mathrm{avg}}}{\log\Lambda}<1
\]
for large enough \(\Lambda\).  Iteration gives a good scale with \(D\to0\).
\end{proof}

\begin{proposition}[Source-doubling scale selection]
\label{prop:source-doubling}
If \(\limsup_{r\downarrow0}P(r)>0\), then for every fixed
\(\Lambda_{\mathrm{par}}>1\) and \(C_{\mathrm{dbl}}>1\) there exist
\(r_n\downarrow0\) and \(p_*>0\) such that
\[
P(r_n)\ge p_*,
\qquad
P(\Lambda_{\mathrm{par}}r_n)\le C_{\mathrm{dbl}}P(r_n).
\]
\end{proposition}

\begin{proof}
Assume otherwise.  Choose \(p_*<\limsup P\).  If every \(r\) with
\(P(r)\ge p_*\) satisfies
\[
P(\Lambda_{\mathrm{par}}r)>C_{\mathrm{dbl}}P(r),
\]
then iteration gives
\[
P(\Lambda_{\mathrm{par}}^m r)>C_{\mathrm{dbl}}^m p_*
\]
while \(\Lambda_{\mathrm{par}}^m r\) remains in a fixed finite parent
cylinder.  This contradicts local source-ledger finiteness on that parent
cylinder as \(m\to\infty\).
\end{proof}

\begin{proposition}[Fixed-dilation promoted-tail payment]
\label{prop:fixed-tail}
Assume promoted-tail slices obey bounded overlap in a parent cylinder:
\[
\sum_{j,k}\mathbf 1_{\mathcal T_{j,k}^{\mathrm{tail}}(r)}
\le
N_{\mathrm{tail}}\mathbf 1_{Q_{\Lambda_{\mathrm{par}}r}(z_0)}+o_r(1).
\]
Along source-doubling radii,
\[
r_n\sum_{j,k}\iint_{\mathcal T_{j,k}^{\mathrm{tail}}(r_n)}
\sigma_+\,dx\,dt
\le
C_{\mathrm{tail}}^{\mathrm{fp}}P(r_n)+o_n(1).
\]
\end{proposition}

\begin{proposition}[Bounded macroscopic-parent tail smallness]
\label{prop:macro-tail}
Assume promoted-tail slices obey bounded overlap in a fixed parent cylinder:
\[
\sum_{j,k}\mathbf 1_{\mathcal T_{j,k}^{\mathrm{tail}}(r_n)}
\le
N_{\mathrm{tail}}\mathbf 1_{Q_{r_0}(z_0)}+o_n(1),
\]
with \(P(r_0)<\infty\).  Then
\[
r_n\sum_{j,k}\iint_{\mathcal T_{j,k}^{\mathrm{tail}}(r_n)}
\sigma_+\,dx\,dt=o_n(1).
\]
\end{proposition}

\begin{proof}
Bounded overlap gives a parent source bound.  Since
\[
r_n\iint_{Q_{r_0}(z_0)}\sigma_+\,dx\,dt
=\frac{r_n}{r_0}P(r_0)\to0,
\]
the fixed macroscopic parent contribution is negligible.
\end{proof}

\clearpage
\section{Adversarial Mechanism Tests}
\label{sec:tests}

The empirical layer is an attack layer.  It checks whether the proposed
payment channels appear in cached turbulent data and adversarial finite-radius
audits.  It
does not prove any continuum inequality.

\subsection*{Glossary of computational labels}

The labels \(c413\), \(c309\), and \(c185\) denote cached promoted-source
candidate families extracted from Johns Hopkins Turbulence Database
diagnostics~\cite{JHTDB2008,JHTDB2007}.  The labels \(D3,D5,D6\) denote
denominator variants used in the tail-denominator ladder audit: \(D5\) is the
local-plus-observed-annular
capacity used by that audit, while \(D3\) and \(D6\) require same-row
moving-tube or velocity-pressure CKN data that were unavailable for those
rows.  The tag ``ARR c185 final81'' denotes the annular-renewal reservoir audit
for the \(c185\) candidate at support size \(81\).

\begin{center}
\scriptsize
\begin{tabular}{@{}p{0.25\linewidth}p{0.42\linewidth}p{0.25\linewidth}@{}}
\toprule
\textbf{Attack or audit} & \textbf{Observed result} & \textbf{Interpretation} \\
\midrule
Near-degenerate same-parent split &
For \(M=1,4,8,16,32,64\), same-parent sub-tubes in cached \(c413/c309\)
families collapse under the physical quotient; captured source remains
\(\eta=1\). &
Artificial label multiplicity is quotient-redundant. \\
\addlinespace
Renewal-cascade cross-parent jitter &
High capture persists, but post-quotient tube growth remains large; available
physical cost stays positive in completed \(c413/c309\) rows. &
The attack stresses physical charge and renewal rather than producing a
low-cost hidden packet. \\
\addlinespace
Tail denominator ladder &
In 15 rows the best available denominator was \(D5\), while \(D3,D6\) were
unresolved.\footnotemark &
Arbitrary tail payment is not automatic; missing denominator families must be
named, not filled by proxy. \\
\addlinespace
ARR \(c185\) final81 attribution &
Renewal-exposure correction moves the ARR ratio from about \(1.022\) to
\(0.967\). &
Renewal exposure is a meaningful accounting channel in this audit. \\
\addlinespace
Coherent residual attribution &
Finite-radius \(c413\) rows route to renewal/deactivation or signed damping,
not to a dangerous coherent positive source. &
Supports the residual-routing taxonomy, but does not prove the coherent row. \\
\bottomrule
\end{tabular}
\end{center}
\footnotetext{The public package includes derived cached rows and notebook
checks, not raw JHTDB dumps or the full extraction pipeline.  The underlying
audit recorded best ratios ranging from about \(1.27\) to \(22.56\).}

\begin{remark}[Evidence fence]
The tests are falsification and mechanism checks.  They do not prove
source-ledger finiteness, K3 quotient-or-charge, K4 forward absorption, K5
reverse matrix closure, K5B high-strain rigidity, or the K6 loop gap.  Their
value is that they refuse to hide missing continuum quantities: unavailable
rows remain marked as unresolved.
\end{remark}

\clearpage
\section{Remaining Theorem Gates}
\label{sec:gates}

The program is useful only if the remaining gates are explicit.  The downstream
technical gates K3 through K7 are isolated in
\texttt{technical\_gates\_k3\_k7.pdf}; the list below is the short dependency
ledger used by this paper.
\begin{enumerate}
\item \textbf{Source-ledger finiteness or truncation.}  Formulate the source
ledger in the intended solution class, or replace \(P\) by a truncation with a
stable limiting argument.

\item \textbf{Primary source tail dichotomy/exhaustion.}  The fixed-dilation
and bounded macroscopic-parent tail cases are controlled; genuinely nonlocal
parents must be source-accounted or routed to charge/leakage.

\item \textbf{K3 quotient-or-charge graph theorem.}  Prove the finite weighted
graph statement \eqref{eq:K3-target-standalone}, including the oscillation
coordinate \(\Phi_{\mathrm{osc}}\).

\item \textbf{K4 forward absorption.}  Prove
\[
P\mathcal F_{\mathrm{phys}}^{\mathrm{final}}
\le C_{\mathrm{forw}}^{\mathrm{final}}D+o.
\]

\item \textbf{K5B high-strain matrix rigidity.}  Prove that source-weighted
high strain on a nonrenewing, nonseparating branch core is paid by shape,
oscillation/twist, renewal, or separation, with the source-weighted path
normalization and the constant \(+1\) routed to \(P\).

\item \textbf{Packed boundary reverse constants.}  The packed boundary row
closes good-scale null-source regularity.  K5 still needs reverse-row constants
for the same packed row.

\item \textbf{K6 spectral and loop-gap certification.}  Supply actual weights
\(w\) and constants with \(q_A(w)<1\) and
\[
C_{\mathrm{forw}}^{\mathrm{final}}C_0<c_{\mathrm{final}}(\eta).
\]

\item \textbf{K7 compatibility under ledger edits.}  Any change to \(D(r)\)
must preserve Proposition~\ref{prop:CKN-bridge}.
\end{enumerate}

\clearpage
\section{Conclusion}
\label{sec:conclusion}

The proposed law is deliberately simple:
\[
\boxed{\mathrm{Amplification}\le\mathrm{Payment}.}
\]
This manuscript defines the payment ledger, states the conditional theorem,
records the partial closures, reports adversarial mechanism tests, and names
the gates that remain.  It is a research program paper, not a completed Clay
proof.  Its usefulness is that it turns a vague blowup question into a finite
set of attackable accounting failures.
The next stage is therefore not to reinterpret the closure theorem, but to
prove or falsify the named analytic gates K3--K7 with explicit constants.

\begin{thebibliography}{99}

\bibitem{Leray1934}
J.~Leray,
Sur le mouvement d'un liquide visqueux emplissant l'espace,
\emph{Acta Math.} \textbf{63} (1934), 193--248.
\href{https://doi.org/10.1007/BF02547354}{doi:10.1007/BF02547354}.

\bibitem{Hopf1951}
E.~Hopf,
\"Uber die Anfangswertaufgabe f\"ur die hydrodynamischen Grundgleichungen,
\emph{Math. Nachr.} \textbf{4} (1951), 213--231.

\bibitem{Prodi1959}
G.~Prodi,
Un teorema di unicit\`a per le equazioni di Navier--Stokes,
\emph{Ann. Mat. Pura Appl.} \textbf{48} (1959), 173--182.

\bibitem{Serrin1962}
J.~Serrin,
On the interior regularity of weak solutions of the Navier--Stokes equations,
\emph{Arch. Rational Mech. Anal.} \textbf{9} (1962), 187--195.
\href{https://doi.org/10.1007/BF00253344}{doi:10.1007/BF00253344}.

\bibitem{CKN1982}
L.~Caffarelli, R.~Kohn, and L.~Nirenberg,
Partial regularity of suitable weak solutions of the Navier--Stokes equations,
\emph{Comm. Pure Appl. Math.} \textbf{35} (1982), 771--831.
\href{https://doi.org/10.1002/cpa.3160350604}{doi:10.1002/cpa.3160350604}.

\bibitem{BKM1984}
J.~T.~Beale, T.~Kato, and A.~Majda,
Remarks on the breakdown of smooth solutions for the 3-D Euler equations,
\emph{Comm. Math. Phys.} \textbf{94} (1984), 61--66.
\href{https://doi.org/10.1007/BF01212349}{doi:10.1007/BF01212349}.

\bibitem{ConstantinFefferman1993}
P.~Constantin and C.~Fefferman,
Direction of vorticity and the problem of global regularity for the
Navier--Stokes equations,
\emph{Indiana Univ. Math. J.} \textbf{42} (1993), 775--789.
\href{https://doi.org/10.1512/iumj.1993.42.42034}{doi:10.1512/iumj.1993.42.42034}.

\bibitem{Lin1998}
F.-H.~Lin,
A new proof of the Caffarelli--Kohn--Nirenberg theorem,
\emph{Comm. Pure Appl. Math.} \textbf{51} (1998), 241--257.
\href{https://doi.org/10.1002/(SICI)1097-0312(199803)51:3\%3C241::AID-CPA2\%3E3.0.CO;2-A}{doi:10.1002/(SICI)1097-0312(199803)51:3\%3C241::AID-CPA2\%3E3.0.CO;2-A}.

\bibitem{ESS2003}
L.~Escauriaza, G.~A.~Seregin, and V.~\v{S}ver\'ak,
\(L^{3,\infty}\)-solutions of Navier--Stokes equations and backward uniqueness,
\emph{Russian Math. Surveys} \textbf{58} (2003), 211--250.
\href{https://doi.org/10.1070/RM2003v058n02ABEH000609}{doi:10.1070/RM2003v058n02ABEH000609}.

\bibitem{JHTDB2008}
Y.~Li et al.,
A public turbulence database cluster and applications to study Lagrangian
evolution of velocity increments in turbulence,
\emph{Journal of Turbulence} \textbf{9} (2008), N31.
\href{https://doi.org/10.1080/14685240802376389}{doi:10.1080/14685240802376389}.

\bibitem{JHTDB2007}
E.~Perlman, R.~Burns, Y.~Li, and C.~Meneveau,
Data exploration of turbulence simulations using a database cluster,
in \emph{SC07: Proceedings of the 2007 ACM/IEEE Conference on
Supercomputing}, 2007.
\href{https://doi.org/10.1145/1362622.1362654}{doi:10.1145/1362622.1362654}.

\bibitem{Hou2009}
T.~Y.~Hou,
Blow-up or no blow-up? A unified computational and analytic approach to
3D incompressible Euler and Navier--Stokes equations,
\emph{Acta Numerica} \textbf{18} (2009), 277--346.
\href{https://doi.org/10.1017/S0962492906420018}{doi:10.1017/S0962492906420018}.

\end{thebibliography}

\end{document}
